\section{Introduction}

\subsection{Background}
Networks of connected devices at the edge, such as IoT gateways, branch
offices and industrial controllers, now carry a growing share of traffic and
an even faster-growing share of attacks~\cite{iot_survey}. Machine-learning
intrusion detection systems (IDS) report high accuracy on benchmark datasets,
using classical ensembles~\cite{pandit2023} or hybrid models that combine
anomaly detection with supervised classification~\cite{bayesh2024}. These
systems share one assumption: all traffic is collected at a central server
for training. At the edge this assumption is costly. Raw flows reveal which
hosts talk to which, and moving them off-site raises privacy and legal
concerns, consumes bandwidth, and adds latency to every model update. Many
sites cannot share raw traffic at all.

\subsection{Motivation}
Federated learning (FL)~\cite{mcmahan2017} addresses part of this problem:
each site trains on its own data and shares only model updates. Federated IDS
have been built with semi-supervised objectives~\cite{aouedi2022fluids} and
with Transformer backbones~\cite{devi2025fatvit}. Four gaps remain. First,
most federated IDS exchange the full model in every round, so communication
grows with model size. Parameter-efficient methods such as LoRA~\cite{hu2022lora}
reduce this, but they are rarely applied to network-flow detection. Second,
attack classes are unevenly spread: an attack may be labelled at only one
site, and it is unclear whether federation actually transfers that knowledge
to the other sites. Third, a single malicious client can replace the global
model, and every site that adopts it blindly inherits the damage. Fourth, most
work stops at classification. A verdict is useful only once it becomes an
action, and an IDS that blocks a host on a single false positive cuts off all
of that host's later legitimate traffic. Agentic approaches to security,
which plan, act and adapt rather than only label~\cite{panwar2025agentic},
are mostly discussed separately from federated detection.

A further, methodological gap concerns evaluation. Federated IDS gains are
usually reported as accuracy or recall at the default decision threshold.
A rise in recall that comes together with a rise in false positives may
reflect only a shift of the operating point. Moving the threshold of the
original model would produce the same rise without any federation. Without
threshold-free metrics, the two cannot be told apart.

\subsection{Problem Statement}
We consider $K$ edge sites, each holding a private, non-IID set of labelled
network flows, with one attack class labelled at only one site. The goal is
to build a system that (i) adapts a shared detector without moving raw flows,
(ii) keeps each model upload small, (iii) remains safe when a client is
malicious, and (iv) turns each verdict into an action that stops attacks
without disrupting legitimate hosts. We call the edge component
\emph{agentic} in a specific, limited sense. It chooses among actions by
expected cost, reverses its own blocks when they prove wrong, learns from
analyst feedback, shares and retracts threat intelligence with its peers,
and verifies a new global model before adopting it. It is therefore more than
a classifier followed by a fixed rule. We also ask whether the federated
detector genuinely learns to separate the new attack class, or only moves
its decision threshold.

\subsection{Contributions}
This paper presents Agentic FedGrid and makes the following contributions:
\begin{enumerate}
\item \textbf{A deployed edge--cloud architecture.} It federates only the
low-rank (LoRA) adapters of a frozen tabular Transformer through serverless
AWS services (API Gateway, Lambda, DynamoDB, S3). Each upload is 10.3\,KB,
compared with 103.7\,KB for the full model, and raw flows never leave the
device. The live backend closed 12{,}610 federated rounds over 25 days,
with straggler handling for partial participation.
\item \textbf{An edge agent whose actions go beyond classification.} It
provides cost-aware action selection, graded and self-expiring blocks,
analyst rollback with learned exceptions, peer threat intelligence with
retraction, and verified adoption of global models. We evaluate it with
both an ideal and an imperfect analyst.
\item \textbf{A threshold-free evaluation of federated adaptation.} For an
attack class withheld from pretraining, the recall gain from federation is
fully explained by a shift of the operating point. This holds even when the
class is spread over three clients or when local training is five times
longer. Default-threshold recall is therefore an unreliable measure of
federated gains.
\item \textbf{A multi-seed evaluation on real cloud hardware.} It covers
LoRA rank versus full-model federation; FedAvg, FedProx and median
aggregation; model poisoning; the local training components; and the agent's
decision layer, together with gradient-boosted and centrally trained
reference models.
\end{enumerate}

\subsection{Paper Organisation}
Section~II reviews related work. Section~III describes the system model and
methodology, and Section~IV the experimental setup. Section~V reports the
results. Section~VI discusses limitations, and Section~VII concludes with
directions for future work.
